\chapter{Results}



This chapter presents the results obtained from the experiments conducted with the proposed framework for explainable brain tumor classification. Various configurations and techniques were evaluated, including StyleGAN-C, CNNs and explainable mechanisms. The results are organized by experiment. Additionally, expert evaluation from medical professionals is included to assess the clinical relevance and interpretability of the proposed approach.


\section{Results: Experiment 1}




\subsection{Compute Budget} 

StyleGAN-C was trained separately on the reduced dataset and on the complete dataset, in both cases using an NVIDIA A100 GPU with 80\,GB of VRAM and a runtime environment with 160\,GB of RAM, which corresponds to the most powerful hardware configuration currently available on Google Colab. The training process lasted approximately \textit{8 hours}, reaching \texttt{kimg=600}. This setup enabled efficient processing of high-resolution MRI slices ($256 \times 256$) and ensured stable convergence of the generative model.

\subsection{Generated Images} Figure~\ref{fig:stylegan_conditional_only1} and \ref{fig:stylegan_conditional_only2} show 9 sample images generated by StyleGAN-C, with 3 images per tumor class. Each row corresponds to one of the three tumor types (glioma, meningioma, and pituitary), demonstrating the model's ability to produce class conditional synthetic MRIs that visually resemble real brain tumor images. It is important to note that this version of the generator is conditioned only on tumor class, not on anatomical plane, which may account for certain visual differences in structure and orientation.

\begin{figure}[H]
\centering

\begin{subfigure}[b]{0.48\textwidth}
    \centering
    \includegraphics[width=\textwidth]{Template-Tesis-dii/img/trunc1complete.png}
    \caption{Complete dataset generator.}
    \label{fig:stylegan_conditional_only1}
\end{subfigure}
\hfill
\begin{subfigure}[b]{0.48\textwidth}
    \centering
    \includegraphics[width=\textwidth]{Template-Tesis-dii/img/trunc1inge.png}
    \caption{Reduced dataset generator.}
    \label{fig:stylegan_conditional_only2}
\end{subfigure}

\caption{Class-conditional samples generated at truncation $\psi = 1.0$. Rows correspond to tumor classes. The left model was trained on the full dataset, while the right was trained on a reduced subset.}
\label{fig:stylegan_conditional_comparison}
\end{figure}


\subsection{Truncation Effect} 

To analyze the effect of truncation on generation diversity and realism, Figure~\ref{fig:stylegan_class_vs_truncation1} and Figure~\ref{fig:stylegan_class_vs_truncation2} show a grid of synthetic images produced by StyleGAN-C. Each row corresponds to a different tumor class (glioma, meningioma, pituitary), while columns vary the truncation parameter from 0.8 to 1.3. Lower truncation values yield more conservative and realistic outputs, whereas higher values produce more diverse but potentially less plausible samples. 

\begin{figure}[H]
    \centering
    \includegraphics[width=0.9\textwidth]{Template-Tesis-dii/img/serialinge.png}
    \caption{Grid of images generated by StyleGAN-C across different tumor classes (rows) and truncation values (columns). Reduced Dataset.}
    \label{fig:stylegan_class_vs_truncation1}
\end{figure}

\begin{figure}[H]
    \centering
    \includegraphics[width=0.9\textwidth]{Template-Tesis-dii/img/serialcomplet.png}
    \caption{Grid of images generated by StyleGAN-C across different tumor classes (rows) and truncation values (columns). Complete Dataset.}
    \label{fig:stylegan_class_vs_truncation2}
\end{figure}


\subsection{Generated Images vs Real Images} Figure~\ref{fig:stylegan_real_vs_fake} compares synthetic and real MR images from the same tumor classes. Each pair of adjacent columns corresponds to a specific class (glioma, meningioma, pituitary), with synthetic samples on the top row and real images on the bottom. 


\begin{figure}[H]
\centering
\includegraphics[width=0.9\textwidth]{Template-Tesis-dii/img/ejemplos/descarga (8).png}
\caption{Comparison between real MR images and synthetic samples. 
The top row shows synthetic images and the bottom row shows real MRIs; pairs of adjacent columns correspond to the three tumor classes.}
\label{fig:stylegan_real_vs_fake}
\end{figure}


\subsection{Latent Inversion} 


Latent inversion was conducted to evaluate the capacity of StyleGAN-C to reconstruct real MR images while enforcing different class embeddings. In Figure~\ref{fig:stylegan_inversion_only}, panel~(a) displays a real pituitary tumor slice from the test set. The following panels~(b)--(d) show reconstructions of the same image using the same optimized latent code but with different class embeddings: glioma, meningioma and pituitary.



\begin{figure}[H]
\centering

\begin{subfigure}[b]{0.33\textwidth}
  \centering
  \includegraphics[width=\textwidth]{Template-Tesis-dii/img/real_pituitary.png}
  \caption{Real Pituitary.}
  \label{fig:stylegan_inversion_1}
\end{subfigure}
\hfill
\begin{subfigure}[b]{0.33\textwidth}
  \centering
  \includegraphics[width=\textwidth]{Template-Tesis-dii/img/cf_to_glioma_f02.png}
  \caption{Reconstruction with the class embedding fixed to \emph{glioma}.}
  \label{fig:stylegan_inversion_2}
\end{subfigure}

\vspace{0.5em}

\begin{subfigure}[b]{0.33\textwidth}
  \centering
  \includegraphics[width=\textwidth]{Template-Tesis-dii/img/cf_to_meningioma_f02.png}
  \caption{Reconstruction with the class embedding fixed to \emph{meningioma}.}
  \label{fig:stylegan_inversion_3}
\end{subfigure}
\hfill
\begin{subfigure}[b]{0.33\textwidth}
  \centering
  \includegraphics[width=\textwidth]{Template-Tesis-dii/img/cf_to_glioma_f02.png}
  \caption{Reconstruction with the class embedding fixed to \emph{pituitary}.}
  \label{fig:stylegan_inversion_4}
\end{subfigure}

\caption{Example of StyleGAN-C inversion on a held-out test MRI. 
Panel~(a) shows the real target image. Panels~(b)--(d) show reconstructions obtained 
by optimizing the latent code while constraining the generator to each tumor class 
(glioma, meningioma, and pituitary, respectively).}
\label{fig:stylegan_inversion_only}
\end{figure}



\subsection{Quantitative Metrics}

Image quality metrics are computed using the conditional generator trained exclusively on the reduced training set, with all generator weights kept fixed during evaluation. At test time, real MRI slices from the held-out testing set are provided as inputs, projected/inverted into the generator's latent space through the inversion module, and re synthesized as class conditioned reconstructions for each target class. Consequently, all evaluated generations correspond to unseen test examples, enabling an unbiased assessment of generalization. Pairwise metrics (PSNR, SSIM, LPIPS) compare each class conditioned reconstruction to its original test image, whereas the set-level metric (FID) measures distributional similarity between real test images and generated images within each class.

\subsubsection{Inversion Images}

\begin{table}[H]
\centering
\caption{Diagonal reconstruction metrics (real class = generated class). PSNR and SSIM are higher-is-better, LPIPS and FID are lower-is-better. PSNR/SSIM/LPIPS are reported as mean $\pm$ std over paired samples, while FID is computed at the set level per class.}
\label{tab:diag_metrics_with_fid}
\begin{tabular}{lcccc}
\toprule
\textbf{Real class} & \textbf{PSNR (dB)} & \textbf{SSIM} & \textbf{LPIPS} & \textbf{FID} \\
\midrule
glioma      & $15.461 \pm 1.856$ & $0.543 \pm 0.093$ & $0.211 \pm 0.071$ & 123.909 \\
meningioma  & $13.908 \pm 1.641$ & $0.474 \pm 0.096$ & $0.280 \pm 0.083$ & 142.130 \\
pituitary   & $13.300 \pm 1.633$ & $0.377 \pm 0.091$ & $0.264 \pm 0.081$ & 168.521 \\
\midrule
total       & $14.427 \pm 1.976$ & $0.478 \pm 0.116$ & $0.245 \pm 0.083$ & 103.958 \\
\bottomrule
\end{tabular}
\end{table}


\begin{table}[ht]
\centering
\caption{Pairwise metrics for all real-to-generated class combinations.}
\label{tab:all_combos_psnr_ssim_lpips}
\begin{tabular}{llccc}
\toprule
\textbf{Real class} & \textbf{Gen class} & \textbf{PSNR (dB)} & \textbf{SSIM} & \textbf{LPIPS} \\
\midrule
glioma     & glioma     & 15.461 & 0.543 & 0.211 \\
glioma     & meningioma & 14.416 & 0.495 & 0.257 \\
glioma     & pituitary  & 13.887 & 0.439 & 0.238 \\
meningioma & glioma     & 14.041 & 0.484 & 0.271 \\
meningioma & meningioma & 13.908 & 0.474 & 0.280 \\
meningioma & pituitary  & 13.204 & 0.429 & 0.306 \\
pituitary  & glioma     & 12.826 & 0.369 & 0.294 \\
pituitary  & meningioma & 12.533 & 0.360 & 0.314 \\
pituitary  & pituitary  & 13.300 & 0.377 & 0.264 \\
\bottomrule
\end{tabular}
\end{table}

\subsubsection{Images reduced data set generator images}

As a result, when comparing the distribution of real training images and the synthetic images generated by StyleGAN-C on the reduced dataset, the best global FID was obtained for truncation $\psi = 1.0$ with a score of 46.08 (Table~\ref{tab:fid_global_truncation}).


\begin{table}[H]
\centering
\caption{Global FID as a function of the truncation parameter $\psi$ and 600 samples per class.}
\label{tab:fid_global_truncation}
\begin{tabular}{lc}
\toprule
\textbf{Truncation $\psi$} & \textbf{FID (global)} \\
\midrule
0.9 & 46.239 \\
1.0 & 46.08 \\
1.1 & 49.192 \\
1.2 & 54.024 \\
1.3 & 58.296 \\
1.4 & 67.509 \\
\bottomrule
\end{tabular}
\end{table}

\clearpage
\section{Results: Experiment 2A}

This experiment presents the performance metrics of the classification model by inference (LPIPS based), in comparison with a CNN baseline. Unlike the other experiments, this setup performs classification at the patient level rather than at the slice level and uses only one slice per patient. This decision was taken due to the high computational cost associated with applying generative inversion and LPIPS  inference to every individual slice. Table~\ref{tab:ensemble_overall} reports the global performance metrics for all three methods, while Table~\ref{tab:ensemble_confusion} shows the corresponding patient-level confusion matrices.

\begin{table}[H]
\centering
\caption{Patient-level performance of the CNN and LPIPS-based inference by generation (68 patients).}
\label{tab:ensemble_overall}
\begin{tabular}{lcccc}
\toprule
Method & Accuracy & Precision (macro) & Recall (macro) & F1-score (macro) \\
\midrule
CNN           & 0.838 & 0.864 & 0.858 & 0.852 \\
LPIPS        & 0.765 & 0.757 & 0.760 & 0.757 \\
\bottomrule
\end{tabular}
\end{table}


\begin{table}[H]
\centering
\caption{Confusion matrices at patient level for the CNN baseline and LPIPS-based inference by generation (rows = true class, columns = predicted class).}
\label{tab:ensemble_confusion}
\begin{tabular}{lccc}
\toprule
 & \multicolumn{3}{c}{Predicted class} \\
\cmidrule(lr){2-4}
True class / Method & glioma & meningioma & pituitary \\
\midrule
\multicolumn{4}{l}{\textbf{CNN (prob-based aggregation)}} \\
glioma     & 21 & 8 & 0 \\
meningioma & 2  & 19 & 0 \\
pituitary  & 0  & 1  & 17 \\
\midrule
\multicolumn{4}{l}{\textbf{LPIPS (normalized inference by generation)}} \\
glioma     & 24 & 2 & 3 \\
meningioma & 6  & 13 & 2 \\
pituitary  & 0  & 3  & 15 \\
\bottomrule
\end{tabular}
\end{table}


\subsection{Counterfactual Videos}

To support visual explainability and clinical interpretability, counterfactual videos were generated for each test patient. These videos show how a real MRI slice would appear if it belonged to a different tumor class, by interpolating the latent space and conditioning the generator on alternative class embeddings. For each patient, the first available slice is selected as input, and a total of 120 frames are generated by linearly interpolating from the original latent code to the class-conditional projection. Each patient is associated with three class conditional videos (one per tumor type) revealing fine or pronounced changes in tumor structure and surrounding tissue throughout the transformation.


\begin{figure}[H]
    \centering
    \includegraphics[width=\textwidth]{Template-Tesis-dii/img/videoexample.png}
    \caption{Counterfactual trajectory for a glioma case (glioma$\rightarrow$glioma). The 120 frames show a latent interpolation that preserves the glioma class while progressively moving towards the corresponding class prototype.}
    \label{fig:counterfactual_120frames}
\end{figure}


\begin{figure}[ht]
    \centering
    \includegraphics[width=\textwidth]{Template-Tesis-dii/img/videoexamplepituitary.png}
    \caption{Counterfactual trajectory for a glioma case (glioma$\rightarrow$pituitary). The 120 frames show a latent interpolation that preserves the glioma class while progressively moving towards the corresponding class prototype.}
    \label{fig:counterfactual_120frames1}
\end{figure}


\clearpage

\section{Results: Experiment 2B}



\begin{enumerate}
    \item \textbf{CNN}: SVM trained on penultimate-layer features from the fine tuned tumor ResNet-50.
    \item \textbf{Encoder}: SVM trained on embeddings produced by the supervised encoder.
    \item \textbf{e4e}: SVM trained on W\textsuperscript{+} latent codes extracted by the e4e encoder pre-trained on StyleGAN2-FFHQ and fine tuned on the tumor dataset.
\end{enumerate}


\subsection{SVM Metrics}

Table~\ref{tab:exp2b_all_metrics} summarizes the global performance on the test set in terms of accuracy and macro-averaged precision, recall and F1-score.

\begin{table}[ht]
    \centering
    \caption{Per-class and global metrics on the test set for the three feature configurations. 
    For class rows, precision/recall/F1-score are per-class; for the total row, they correspond to macro-averaged precision, recall and F1-score.}
    \label{tab:exp2b_all_metrics}
    \begin{tabular}{llcccc}
        \toprule
        Setting        & Class       & Precision & Recall  & F1-score & Accuracy \\
        \midrule
        % ============== CNN-only ==============
        \multirow{4}{*}{CNN}
                       & glioma      & 0.961     & 0.939   & 0.950    & --       \\
                       & meningioma  & 0.847     & 0.887   & 0.866    & --       \\
                       & pituitary   & 0.967     & 0.970   & 0.969    & --       \\
        \cmidrule(lr){2-6}
                       & \textbf{total}     
                                    & \textbf{0.925} 
                                              & \textbf{0.932} 
                                                      & \textbf{0.928} 
                                                              & \textbf{0.9374}   \\
        \midrule
        % ============== Encoder-only ==========
        \multirow{4}{*}{Encoder}
                       & glioma      & 0.970     & 0.948   & 0.959    & --       \\
                       & meningioma  & 0.851     & 0.912   & 0.881    & --       \\
                       & pituitary   & 0.985     & 0.970   & 0.978    & --       \\
        \cmidrule(lr){2-6}
                       & \textbf{total}     
                                    & \textbf{0.935} 
                                              & \textbf{0.944} 
                                                      & \textbf{0.939} 
                                                              & \textbf{0.9473}   \\ 
        \midrule
        % ============== E4E-only ==============
        \multirow{4}{*}{e4e}
                       & glioma      & 0.864     & 0.852   & 0.858    & --       \\
                       & meningioma  & 0.643     & 0.696   & 0.668    & --       \\
                       & pituitary   & 0.931     & 0.897   & 0.914    & --       \\
        \cmidrule(lr){2-6}
                       & \textbf{total}     
                                    & \textbf{0.813} 
                                              & \textbf{0.815} 
                                                      & \textbf{0.813} 
                                                              & \textbf{0.8321}   \\
        \bottomrule
    \end{tabular}
\end{table}


The supervised encoder achieves the best overall performance, with an accuracy of 0.9473 and a macro F1-score of 0.9390, slightly outperforming the CNN features (accuracy = 0.9374, macro F1 = 0.9280). In contrast, the e4e configuration yields substantially lower performance (accuracy = 0.8321, macro F1 = 0.8130).


\subsection{Confusion matrices}

Table~\ref{tab:cm_all_models} reports the confusion matrices for the three feature configurations, with rows corresponding to the true class and columns to the predicted class.

\begin{table}[ht]
    \centering
    \caption{Confusion matrices for the three feature configurations on the test set (rows = true class, columns = predicted class).}
    \label{tab:cm_all_models}
    \begin{tabular}{llccc}
        \toprule
        Setting        & True class & glioma & meningioma & pituitary \\
        \midrule
        % ============== CNN-only ==============
        \multirow{3}{*}{CNN}
                       & glioma      & 419   & 24        & 3         \\
                       & meningioma  & 16    & 172       & 6         \\
                       & pituitary   & 1     & 7         & 263       \\
        \midrule
        % ============== Encoder-only ==========
        \multirow{3}{*}{Encoder}
                       & glioma      & 423   & 23        & 0         \\
                       & meningioma  & 13    & 177       & 4         \\
                       & pituitary   & 0     & 8         & 263       \\
        \midrule
        % ============== E4E-only ==============
        \multirow{3}{*}{e4e}
                       & glioma      & 380   & 58        & 8         \\
                       & meningioma  & 49    & 135       & 10        \\
                       & pituitary   & 11    & 17        & 243       \\
        \bottomrule
    \end{tabular}
\end{table}

Overall, the CNN and encoder configurations exhibit similar error patterns, with most mistakes occurring between glioma and meningioma and very few pituitary tumors being confused with the other classes. The encoder model slightly reduces the number of misclassifications for glioma and meningioma compared with the CNN baseline, in line with its higher macro F1-score. In contrast, the e4e configuration produces substantially more errors for all three classes, particularly an increased confusion of meningioma and pituitary with glioma, which is consistent with its lower overall accuracy and macro F1-score.
\subsection{PCA Analysis}

Figure~\ref{fig:pca_exp2b} presents the three-dimensional PCA projections of the feature spaces obtained from the three configurations, separately for the training and test splits. Each point corresponds to an MRI slice and is color coded according to the tumor class, allowing visual inspection of how the different representations distribute the classes in a low dimensional space. Table~\ref{tab:pca_variance} reports the proportion of variance explained by the first three principal components for each feature space and split. Together, the figure and the table summarize the global geometric structure of the learned embeddings and how much of their variability is captured by the leading PCA directions.


\begin{table}[H]
    \centering
    \caption{Explained variance ratio of the first three principal components
             for each feature configuration. The \emph{Shape} column indicates
             the number of samples and feature dimensionality used to fit PCA.}
    \label{tab:pca_variance}
    \begin{tabular}{llccccc}
        \toprule
        Features      & Split & Shape              & PC1    & PC2    & PC3    & SUM(PC1--3) \\
        \midrule
        \multirow{2}{*}{Encoder} 
                      & train & \(2153 \times 512\)  & 0.5893 & 0.4074 & 0.0007 & 0.9974 \\
                      & test  & \(911 \times 512\)   & 0.5809 & 0.4128 & 0.0016 & 0.9952 \\
        \midrule
        \multirow{2}{*}{CNN} 
                      & train & \(2153 \times 2048\) & 0.3698 & 0.3018 & 0.0501 & 0.7216 \\
                      & test  & \(911 \times 2048\)  & 0.3788 & 0.2677 & 0.0454 & 0.6918 \\
        \midrule
        \multirow{2}{*}{e4e} 
                      & train & \(2153 \times 7168\) & 0.3347 & 0.1946 & 0.1337 & 0.6630 \\
                      & test  & \(911 \times 7168\)  & 0.3598 & 0.1780 & 0.1479 & 0.6857 \\
        \bottomrule
    \end{tabular}
\end{table}



\begin{figure}[h]
    \centering

    % ===== Row 1 =====
    \begin{subfigure}[b]{0.4\textwidth}
        \centering
        \includegraphics[width=\textwidth]{Template-Tesis-dii/img/ejemplos/pca/pcwtrain.png}
        \caption{Encoder (train)}
        \label{fig:pca_w_train}
    \end{subfigure}
    \hfill
    \begin{subfigure}[b]{0.4\textwidth}
        \centering
        \includegraphics[width=\textwidth]{Template-Tesis-dii/img/ejemplos/pca/pcawtest.png}
        \caption{Encoder (test)}
        \label{fig:pca_w_test}
    \end{subfigure}

    \vspace{0.5em}

    % ===== Row 2 =====
    \begin{subfigure}[b]{0.4\textwidth}
        \centering
        \includegraphics[width=\textwidth]{Template-Tesis-dii/img/ejemplos/pca/pcacnntrain.png}
        \caption{CNN (train)}
        \label{fig:pca_cnn_train}
    \end{subfigure}
    \hfill
    \begin{subfigure}[b]{0.4\textwidth}
        \centering
        \includegraphics[width=\textwidth]{Template-Tesis-dii/img/ejemplos/pca/pcacnntest.png}
        \caption{CNN (test)}
        \label{fig:pca_cnn_test}
    \end{subfigure}

    \vspace{0.5em}

    % ===== Row 3 =====
    \begin{subfigure}[b]{0.4\textwidth}
        \centering
        \includegraphics[width=\textwidth]{Template-Tesis-dii/img/e4epcatrain.png}
        \caption{e4e (train)}
        \label{fig:pca_fusion_train}
    \end{subfigure}
    \hfill
    \begin{subfigure}[b]{0.4\textwidth}
        \centering
        \includegraphics[width=\textwidth]{Template-Tesis-dii/img/e4epcatest.png}
        \caption{e4e (test)}
        \label{fig:pca_fusion_test}
    \end{subfigure}

    \caption{PCA projections of the encoder, CNN and e4e feature spaces on the training (left column) and test (right column) splits. Each point corresponds to an MRI slice, color coded by tumor class.}
    \label{fig:pca_exp2b}
\end{figure}

\clearpage


\section{Results: Experiment 2C}
\label{sec:results-augmentation}



Experiment 2C evaluates whether class conditional synthetic augmentation with StyleGAN-C improves tumor classification performance under patient level data splits. The study compares real only training against augmented synthetic training across multiple folds, using both the complete and reduced datasets. Additionally, ablation studies explore the influence of truncation levels (\(\psi\)), augmentation strength, CAS, and oversampling baselines.


\begin{table}[ht]
\centering
\caption{CNN Training and augmentation hyperparameters used in all experiments.}
\label{tab:hyperparams}
\begin{tabular}{ll}
\toprule
Hyperparameter                    & Value \\
\midrule
Backbone architecture            & ResNet-50 (ImageNet-pretrained) \\
Input resolution                 & $256\times256$\\
Learning rate                    & $1\times10^{-4}$ \\
Weight decay                     & $1\times10^{-4}$ \\
Batch size                       & 32 \\
Training epochs per fold         & 20 \\
Cross-validation folds $K$       & $K=3$ and $K=5$  \\
Split strategy                   & StratifiedGroupKFold by patient ID \\
GAN backbone                     & Class-conditional StyleGAN2-ADA \\
GAN truncation $\psi$           & $\psi \in \{0.8,0.9,1.0,1.1,1.2,1.3, 1.4, 1.5\}$  \\
Random seed                      & $122$ \\
\bottomrule
\end{tabular}
\end{table}

\paragraph{Dataset comparison}
Parallel experiments were conducted using both the complete and reduced datasets to evaluate the robustness and generalization capabilities of the classification model under varying levels of data availability. Both configurations followed the same experimental protocol, including training with real samples and synthetic augmentation. This setup enables a direct comparison of model performance depending on the amount of real data available and highlights the relative benefits of incorporating generated samples.

\subsection{Complete Dataset}
\label{subsec:aug-primary}

The complete dataset setting uses all available MRI slices from the dataset, providing the classifier with the maximum amount of real data.




\subsubsection{Cross Validation Summary}

The tables present the performance metrics of the model across different cross-validation folds. They compare the results for real images training and GAN augmented training with various truncation values, showing the mean and standard deviation for each metric.

\begin{table}[H]
\centering
\small
\caption{TEST macro metrics for 3-fold experiments: real-only vs.\ GAN-balanced ($\psi=1.4$). Values are mean $\pm$ std across folds.}
\label{tab:3k_test_macro}
\begin{tabular}{lcccc}
\toprule
Setting & Accuracy & Precision (macro) & Recall (macro) & F1 (macro) \\
\midrule
Real-only        & $0.930 \pm 0.014$ & $0.920 \pm 0.016$ & $0.927 \pm 0.006$ & $0.922 \pm 0.013$ \\
GAN ($\psi=1.4$) & $0.945 \pm 0.003$ & $0.940 \pm 0.002$ & $0.938 \pm 0.003$ & $0.939 \pm 0.002$ \\
\bottomrule
\end{tabular}
\end{table}
\begin{table}[H]
\centering
\caption{TEST performance for 5-fold runs using macro metrics (mean $\pm$ standard deviation across folds).}
\label{tab:test_5fold_macro_real_vs_gan}
\begin{tabular}{lcccc}
\toprule
Setting & Accuracy & Precision (macro) & Recall (macro) & F1 (macro) \\
\midrule
Real-only         & $0.941 \pm 0.011$ & $0.934 \pm 0.014$ & $0.935 \pm 0.013$ & $0.934 \pm 0.013$ \\
GAN ($\psi=1.3$)  & $0.941 \pm 0.010$ & $0.932 \pm 0.014$ & $0.937 \pm 0.009$  & $0.934 \pm 0.011$ \\
GAN ($\psi=1.4$)  & $0.948 \pm 0.002$ & $0.941 \pm 0.004$ & $0.942 \pm 0.006$ & $0.941 \pm 0.003$ \\
GAN ($\psi=1.5$)  & $0.937 \pm 0.003$ & $0.926 \pm 0.005$ & $0.933 \pm 0.005$ & $0.929 \pm 0.004$ \\
\bottomrule
\end{tabular}
\end{table}



\subsubsection{Ensemble Evaluation}

The tables present the performance metrics of the ensemble model, which aggregates predictions across multiple cross-validation folds to improve performance and robustness.


\begin{table}[H]
    \centering
    \caption{5 fold Overall macro-averaged performance for the CNN ensemble with different truncation values $\psi$.}
    \label{tab:macro_overall_truncation}
    \begin{tabular}{lcccc}
        \toprule
        Setting        & Accuracy & Precision (macro) & Recall (macro) & F1 (macro) \\
        \midrule
        Real ensemble  & 0.947 & 0.940 & 0.943 & 0.941 \\
        $\psi = 1.3$   & 0.957 & 0.952 & 0.954 & 0.953 \\
        $\psi = 1.4$   & 0.956 & 0.953 & 0.950 & 0.951 \\
        $\psi = 1.5$   & 0.952 & 0.943 & 0.949 & 0.946 \\
        \bottomrule
    \end{tabular}
\end{table}

\begin{table}[H]
\centering
\caption{Per-class test performance of the 3-fold ensemble for real-only vs.\ GAN-augmented training.}
\label{tab:test_ensemble_per_class_3fold_real_vs_gan}
\begin{tabular}{lcccccc}
\toprule
Setting      & Class       & Precision & Recall & F1-score & Accuracy & Test samples \\
\midrule
Real-only    & Glioma      & $0.965$ & $0.930$ & $0.947$ & --      & $446$ \\
Real-only    & Meningioma  & $0.850$ & $0.907$ & $0.878$ & --      & $194$ \\
Real-only    & Pituitary   & $0.974$ & $0.985$ & $0.980$ & --      & $271$ \\
Real-only    & \textbf{Total}
             & $\mathbf{0.930}$ & $\mathbf{0.941}$ & $\mathbf{0.935}$ & $\mathbf{0.942}$ & $\mathbf{911}$ \\
\midrule
GAN ($\psi=1.4$) & Glioma     & $0.953$ & $0.962$ & $0.958$ & --      & $446$ \\
GAN ($\psi=1.4$) & Meningioma & $0.906$ & $0.897$ & $0.902$ & --      & $194$ \\
GAN ($\psi=1.4$) & Pituitary  & $0.970$ & $0.963$ & $0.967$ & --      & $271$ \\
GAN ($\psi=1.4$) & \textbf{Total}
             & $\mathbf{0.943}$ & $\mathbf{0.941}$ & $\mathbf{0.942}$ & $\mathbf{0.948}$ & $\mathbf{911}$ \\
\bottomrule
\end{tabular}
\end{table}

\begin{table}[H]
\centering
\small
\caption{Per-class test performance of the 5-fold ensemble for real-only vs.\ GAN-balanced ($\psi=1.4$).}
\label{tab:ensemble5k_perclass_real_vs_gan14}
\begin{tabular}{lcccccc}
\toprule
Setting & Class & Precision & Recall & F1-score & Accuracy & Test samples \\
\midrule
Real-only        & Glioma     & 0.955 & 0.953 & 0.954 & --     & 446 \\
Real-only        & Meningioma & 0.881 & 0.912 & 0.896 & --     & 194 \\
Real-only        & Pituitary  & 0.985 & 0.963 & 0.974 & --     & 271 \\
Real-only        & \textbf{Total} 
                 & \textbf{0.940} & \textbf{0.943} & \textbf{0.941} & \textbf{0.947} & \textbf{911} \\
\midrule
GAN ($\psi=1.4$) & Glioma     & 0.953 & 0.964 & 0.959 & --     & 446 \\
GAN ($\psi=1.4$) & Meningioma & 0.912 & 0.907 & 0.910 & --     & 194 \\
GAN ($\psi=1.4$) & Pituitary  & 0.993 & 0.978 & 0.985 & --     & 271 \\
GAN ($\psi=1.4$) & \textbf{Total} 
                 & \textbf{0.953} & \textbf{0.950} & \textbf{0.951} & \textbf{0.956} & \textbf{911} \\
\bottomrule
\end{tabular}
\end{table}

\subsection{Reduced Dataset}

The reduced dataset is filtered to exclude MRI slices with poor quality or images where the tumor is not clearly visible, ensuring that only relevant and high-quality data is used for training.

\subsubsection{Baseline}

As a reference configuration, the classifier was trained on all available internal slices and evaluated once on the test set, yielding a single baseline operating point.





\begin{table}[ht]
\centering
\caption{Full-train test performance (macro metrics).}
\begin{tabular}{lccccc}
\toprule
\textbf{Setting} & \textbf{N} & \textbf{Accuracy} & \textbf{Precision (macro)} & \textbf{Recall (macro)} & \textbf{F1 (macro)} \\
\midrule
Full train 
& 702 
& 0.934 
& 0.922 
& 0.933 
& 0.927 \\
\bottomrule
\end{tabular}
\label{tab:full_train_macro}
\end{table}


\subsubsection{Cross Validation Summary}

\begin{table}[H]
\centering
\caption{Test performance for 5-fold experiments with real-only training and synthetic augmentation using different truncation $\psi$.}
\label{tab:test_truncation}
\begin{tabular}{lcccc}
\toprule
Setting        & Accuracy           & Precision (macro)    & Recall (macro)       & F1 (macro)          \\
\midrule
Real only      & $0.948 \pm 0.009$  & $0.934 \pm 0.010$    & $0.945 \pm 0.012$    & $0.939 \pm 0.011$   \\
$\psi = 0.8$   & $0.950 \pm 0.006$  & $0.938 \pm 0.007$    & $0.946 \pm 0.006$    & $0.942 \pm 0.007$   \\
$\psi = 1.1$   & $0.953 \pm 0.009$  & $0.941 \pm 0.011$    & $0.950 \pm 0.011$    & $0.945 \pm 0.011$   \\
$\psi = 1.2$   & $0.952 \pm 0.008$  & $0.940 \pm 0.012$    & $0.948 \pm 0.006$    & $0.944 \pm 0.009$   \\
$\psi = 1.3$   & $0.954 \pm 0.007$  & $0.944 \pm 0.007$    & $0.949 \pm 0.011$    & $0.946 \pm 0.009$   \\
$\psi = 1.4$   & $0.950 \pm 0.005$  & $0.940 \pm 0.006$    & $0.944 \pm 0.006$    & $0.942 \pm 0.006$   \\
\bottomrule
\end{tabular}
\end{table}

\subsubsection{Ensemble Evaluation}

\begin{table}[H]
\centering
\caption{Test performance of the 5-fold ensemble for real-only vs.\ GAN training with truncation $\psi = 1.3$.}
\label{tab:test_ensemble_per_class_5fold}
\begin{tabular}{ccccccc}
\toprule
Setting      & Class      & Precision & Recall & F1-score & Accuracy & Test samples \\
\midrule
Real-only    & Glioma     & $0.964$ & $0.952$ & $0.958$ & --      & $312$ \\
Real-only    & Meningioma & $0.884$ & $0.915$ & $0.899$ & --      & $141$ \\
Real-only    & Pituitary  & $0.980$ & $0.976$ & $0.978$ & --      & $249$ \\
Real-only    & \textbf{Total} 
             & $\mathbf{0.943}$ & $\mathbf{0.948}$ & $\mathbf{0.945}$ & $\mathbf{0.953}$ & $\mathbf{702}$ \\
\midrule
$\psi = 1.3$ & Glioma     & $0.958$ & $0.962$ & $0.960$ & --      & $312$ \\
$\psi = 1.3$ & Meningioma & $0.909$ & $0.922$ & $0.915$ & --      & $141$ \\
$\psi = 1.3$ & Pituitary  & $0.992$ & $0.980$ & $0.986$ & --      & $249$ \\
$\psi = 1.3$ & \textbf{Total} 
             & $\mathbf{0.953}$ & $\mathbf{0.954}$ & $\mathbf{0.954}$ & $\mathbf{0.960}$ & $\mathbf{702}$ \\
\bottomrule
\end{tabular}
\end{table}

\begin{table}[ht]
    \centering
    \caption{Overall macro-averaged performance for the CNN ensemble with different truncation values $\psi$.}
    \label{tab:macro_overall_truncation_inttest}
    \begin{tabular}{lcccc}
        \toprule
        Setting        & Accuracy & Precision (macro) & Recall (macro) & F1 (macro) \\
        \midrule
        Real ensemble  & 0.953 & 0.939 & 0.952 & 0.945 \\
        $\psi = 0.8$   & 0.954 & 0.942 & 0.950 & 0.946 \\
        $\psi = 1.1$   & 0.956 & 0.943 & 0.954 & 0.948 \\
        $\psi = 1.2$   & 0.960 & 0.949 & 0.957 & 0.953 \\
        $\psi = 1.3$   & 0.960 & 0.953 & 0.954 & 0.954 \\
        $\psi = 1.4$   & 0.953 & 0.943 & 0.948 & 0.945 \\
        \bottomrule
    \end{tabular}
\end{table}

\subsection{Ablation}

Different ablations of data augmentation were performed to determine whether the improvements in metrics were due to class balancing or the quality of the synthetic images. To investigate this, ablation studies were conducted with class balancing using duplicated images, as well as weak and strong geometric variations. Additionally, the CAS metric was included to quantify the quality of the generated images and assess how useful these images are for training the model.

\subsubsection{Ablation 1: CAS}


\begin{table}[H]
\centering
\caption{CAS macro metrics on synthetic-only test for different $N$ and truncation $\psi$ (rounded to three decimals).}
\label{tab:cas_synth_metrics_boxed}
\begin{tabular}{cccccc}
\toprule

$N$ & $\psi$ & Accuracy & Precision (macro) & Recall (macro) & F1 (macro) \\
\midrule
400 & 0.8 & 0.905 & 0.887 & 0.908 & 0.894 \\
400 & 1.0 & 0.923 & 0.906 & 0.923 & 0.912 \\
400 & 1.2 & 0.913 & 0.897 & 0.911 & 0.902 \\
400 & 1.4 & 0.909 & 0.892 & 0.889 & 0.890 \\
\midrule
600 & 0.8 & 0.895 & 0.875 & 0.892 & 0.880 \\
600 & 1.0 & 0.899 & 0.878 & 0.889 & 0.883 \\
600 & 1.2 & 0.906 & 0.886 & 0.891 & 0.888 \\
600 & 1.4 & 0.906 & 0.888 & 0.886 & 0.887 \\
\bottomrule
\end{tabular}
\end{table}


\begin{figure}[H]
  \centering
  \includegraphics[width=0.75\textwidth]{Template-Tesis-dii/img/cas.png}
  \caption{Macro F1 on synthetic-only CAS vs.\ truncation $\psi$  (higher is better).}
  \label{fig:cas_macro_f1_truncation}
\end{figure}





\subsubsection{Ablation 2: Oversample Augmentation}

\begin{table}[H]
\centering
\caption{Slice-level performance of the real-only oversampling ensemble baseline on the 5k test set ($n = 702$).}
\begin{tabular}{lc}
\toprule
Metric & Value \\
\midrule
Accuracy             & 0.934 \\
Precision (macro)    & 0.918 \\
Recall (macro)       & 0.925 \\
F1-score (macro)     & 0.921 \\
Precision (weighted) & 0.937 \\
Recall (weighted)    & 0.934 \\
F1-score (weighted)  & 0.935 \\
\bottomrule
\end{tabular}
\label{tab:exp2c_b1_oversample}
\end{table}



\subsubsection{Ablation 3: Weak vs Strong Augmentation}


\begin{table}[H]
\centering
\caption{Comparison between weak and strong augmentation regimes (Test, $n = 702$ slices). Strong augmentation consistently outperforms the weak baseline.}
\begin{tabular}{lccc}
\toprule
Metric & Weak aug & Strong aug & $\Delta$ (Strong -- Weak) \\
\midrule
Accuracy             & 0.932 & 0.959 & +0.027 \\
Precision (macro)    & 0.915 & 0.951 & +0.037 \\
Recall (macro)       & 0.921 & 0.953 & +0.031 \\
F1-score (macro)     & 0.918 & 0.952 & +0.034 \\
Precision (weighted) & 0.934 & 0.959 & +0.026 \\
Recall (weighted)    & 0.932 & 0.959 & +0.027 \\
F1-score (weighted)  & 0.932 & 0.959 & +0.026 \\
\bottomrule
\end{tabular}
\label{tab:exp2c_b2_weak_vs_strong}
\end{table}










\subsection{Statistical Tests}

Statistical tests were performed to assess the significance of model performance across different datasets. The McNemar test compares slice level performance between single run and ensemble models, as well as real only vs. GAN-augmented ensembles. The NLL comparison between the reduced and complete datasets reveals the effect of dataset quality and size on the model's confidence and generalization, with significant differences observed for certain tumor classes.

\begin{table}[ht]
\centering
\caption{Slice-level performance for McNemar test P-value}
\begin{tabular}{lcc}
\toprule
\textbf{Comparison} & \textbf{Complete Dataset} & \textbf{Reduced Dataset} \\
\midrule
Single Run GAN vs Ensemble GAN & 0.000308 & 0.002967 \\
Ensemble Real vs Ensemble GAN & 0.080 & 0.332 \\
\bottomrule
\end{tabular}
\label{tab:mcnemar_comparison_side_by_side}
\end{table}


\begin{table}[ht]
\centering
\caption{Global and per-class slice-level NLL comparison for the reduced and complete datasets (A: real-only, B: synthetic) N(=\#Reducedimages/\#Completeimages).}
\begin{tabular}{lccc}
\toprule
\textbf{Comparison} & \textbf{Reduced Dataset} & \textbf{Complete Dataset} \\
\midrule
\multicolumn{3}{c}{\textbf{Global NLL Comparison (A real-only, B synthetic)}} \\
\midrule
Mean NLL (A) & 0.105 & 0.122 \\
Mean NLL (B) & 0.091 & 0.128 \\
Diff (A - B) & 0.014 & -0.006 \\
\textbf{p-value (two-sided)} & \textbf{0.0065} & \textbf{0.5075} \\
\midrule
\multicolumn{3}{c}{\textbf{Per-Class NLL Comparison}} \\
\midrule
\textbf{Class: Glioma (N=312/446)} & \textbf{Reduced Dataset} & \textbf{Complete Dataset} \\
Mean NLL (A) & 0.124 & 0.086 \\
Mean NLL (B) & 0.086 & 0.095 \\
Diff (A - B) & 0.038 & -0.008 \\
\textbf{p-value (two-sided)} & \textbf{$5.00 \times 10^{-6}$} & \textbf{0.3095} \\
\midrule
\textbf{Class: Meningioma (N=141/194)} & \textbf{Reduced Dataset} & \textbf{Complete Dataset} \\
Mean NLL (A) & 0.160 & 0.253 \\
Mean NLL (B) & 0.175 & 0.307 \\
Diff (A - B) & -0.014 & -0.054 \\
\textbf{p-value (two-sided)} & \textbf{0.4564} & \textbf{0.1374} \\
\midrule
\textbf{Class: Pituitary (N=249/271)} & \textbf{Reduced Dataset} & \textbf{Complete Dataset} \\
Mean NLL (A) & 0.051 & 0.086 \\
Mean NLL (B) & 0.049 & 0.054 \\
Diff (A - B) & 0.002 & 0.032 \\
\textbf{p-value (two-sided)} & \textbf{0.7125} & \textbf{$5.00 \times 10^{-6}$} \\
\bottomrule
\end{tabular}
\label{tab:nll_comparison}
\end{table}

\begin{table}[ht]
\centering
\caption{Global and per-class independent-samples t-test on slice-level NLL for the reduced and complete datasets (Model A: strong oversampling, Model B: synthetic images). $N_A/N_B$ denotes the number of slices for Model A / Model B.}
\begin{tabular}{lccc}
\toprule
\textbf{Comparison} & \textbf{Reduced Dataset} & \textbf{Complete Dataset} \\
\midrule
\multicolumn{3}{c}{\textbf{Global t-test on slice-level NLL (A vs B)}} \\
\midrule
$N_A / N_B$           & 702 / 702       & 702 / 911      \\
Mean NLL (A)          & 0.105           & 0.122          \\
Mean NLL (B)          & 0.091           & 0.128          \\
Diff (A -- B)         & 0.0140          & -0.0060        \\
$t$-statistic         & 2.596           & 0.549          \\
\textbf{p-value (two-sided)} & \textbf{0.0095} & \textbf{0.5829} \\
\midrule
\multicolumn{3}{c}{\textbf{Per-class t-test on slice-level NLL}} \\
\midrule
\textbf{Class: Glioma} & \textbf{Reduced Dataset} & \textbf{Complete Dataset} \\
$N_A / N_B$           & 312 / 312       & 312 / 446      \\
Mean NLL (A)          & 0.0826          & 0.0826         \\
Mean NLL (B)          & 0.0862          & 0.0946         \\
Diff (A -- B)         & -0.0036         & -0.0120        \\
$t$-statistic         & -0.158          & -0.523         \\
\textbf{p-value (two-sided)} & \textbf{0.8748} & \textbf{0.6014} \\
\midrule
\textbf{Class: Meningioma} & \textbf{Reduced Dataset} & \textbf{Complete Dataset} \\
$N_A / N_B$           & 141 / 141       & 141 / 194      \\
Mean NLL (A)          & 0.2462          & 0.2462         \\
Mean NLL (B)          & 0.1745          & 0.3066         \\
Diff (A -- B)         & 0.0717          & -0.0604        \\
$t$-statistic         & 1.082           & -0.699         \\
\textbf{p-value (two-sided)} & \textbf{0.2804} & \textbf{0.4850} \\
\midrule
\textbf{Class: Pituitary} & \textbf{Reduced Dataset} & \textbf{Complete Dataset} \\
$N_A / N_B$           & 249 / 249       & 249 / 271      \\
Mean NLL (A)          & 0.1538          & 0.1538         \\
Mean NLL (B)          & 0.0493          & 0.0542         \\
Diff (A -- B)         & 0.1045          & 0.0996         \\
$t$-statistic         & 4.118           & 3.446          \\
\textbf{p-value (two-sided)} & \textbf{$4.73 \times 10^{-5}$} & \textbf{0.0006} \\
\bottomrule
\end{tabular}
\label{tab:nll_ttest_reduced_complete}
\end{table}

From a probabilistic perspective, the most relevant improvement brought by GAN-based augmentation is observed in the negative log-likelihood (NLL) rather than in raw accuracy. Although the accuracy gains over the only real images ensemble are numerically modest, the GAN-balanced models consistently achieve lower NLL, particularly for the meningioma class. This indicates that the classifier becomes more confident and better calibrated in its probability estimates, assigning higher probability mass to the correct class while avoiding over confident incorrect predictions.


\clearpage




\section{Results: Experiment 3}

\subsection{Grad-Cam}


\begin{table}[ht]
    \centering
    \caption{Grad-CAM visualizations for two representative cases per class using the CNN model.}
    \label{tab:gradcam_cnn}
    \begin{tabular}{lcc}
        \toprule
        Class & Grad-CAM example 1 & Grad-CAM example 2 \\
        \midrule
        
        Glioma &
        \includegraphics[width=0.22\textwidth]{Template-Tesis-dii/img/ejemplos/xai/glioma_true__glioma_pred__MR036496B_glioma_0002.png} &
        \includegraphics[width=0.22\textwidth]{Template-Tesis-dii/img/ejemplos/xai/glioma_true__meningioma_pred__107747_glioma_0012.png} \\
        
        Meningioma &
        \includegraphics[width=0.22\textwidth]{Template-Tesis-dii/img/ejemplos/xai/meningioma_true__meningioma_pred__97607_meningioma_0002.png} &
        \includegraphics[width=0.22\textwidth]{Template-Tesis-dii/img/ejemplos/xai/meningioma_true__meningioma_pred__113435_meningioma_0017.png} \\
        
        Pituitary &
        \includegraphics[width=0.22\textwidth]{Template-Tesis-dii/img/ejemplos/xai/pituitary_true__pituitary_pred__112649_pituitary_0006.png} &
        \includegraphics[width=0.22\textwidth]{Template-Tesis-dii/img/ejemplos/xai/pituitary_true__pituitary_pred__112847_pituitary_0002.png} \\
        
        \bottomrule
    \end{tabular}
\end{table}





\section{Quantitative Evaluation}

The quantitative analysis of the HITL experiment is structured in two parts. First, the neuroradiologist's responses are used to assess the diagnostic performance and perceived quality of the images, comparing real MRIs against synthetic samples generated by the proposed model. This part focuses on accuracy, subjective confidence in the predicted tumor class, and perceived diagnostic utility across tumor classes and image types. Second, the study evaluates the behavior of the explainability component, analyzing whether Grad-CAM heat maps correctly highlight tumor regions and how this affects the diagnostic accuracy, confidence and perceived usefulness of the explanations. In both parts, accuracy, neuroradiologist confidence ratings, and diagnostic utility scores are used as the main quantitative endpoints, providing a complementary view of how generated images and XAI overlays support or limit expert decision making.


\subsection{Evaluation of Generated Images}

In this part of the experiment, the quality and diagnostic value of the generated images were evaluated through a blinded reading study. A mixed set of real MRIs and synthetic slices produced by the proposed generator was assembled and presented to the neuroradiologist in random order. The expert was not informed of the ground truth tumor class, nor whether each image was real or synthetic, and was asked to provide an independent diagnosis for every case.


\begin{table}[H]
\centering
\caption{Neuroradiologist performance by image type and tumor class (all cases; unanswered counted as errors).}
\label{tab:hitl_all_cases}
\begin{tabular}{lrrrr}
\toprule
Subset & $n$ & Accuracy & Mean confidence & Mean utility \\
\midrule
All images         & 33 & 0.73 & 2.91 & 2.88 \\
\midrule
Real (all)         & 15 & 1.00 & 3.73 & 3.80 \\
\quad Real (glioma)      &  5 & 1.00 & 2.80 & 2.60 \\
\quad Real (meningioma)  &  5 & 1.00 & 3.40 & 3.80 \\
\quad Real (pituitary)   &  5 & 1.00 & 5.00 & 5.00 \\
\midrule
Fake (all)         & 18 & 0.50 & 2.22 & 2.11 \\
\quad Fake (glioma)      &  6 & 0.50 & 1.67 & 1.67 \\
\quad Fake (meningioma)  &  6 & 0.17 & 1.67 & 1.33 \\
\quad Fake (pituitary)   &  6 & 0.83 & 3.33 & 3.33 \\
\bottomrule
\end{tabular}
\end{table}


\begin{table}[H]
\centering
\caption{Neuroradiologist performance by image type and tumor class (answered cases only; rows with non-missing predicted tumor).}
\label{tab:hitl_answered_only}
\begin{tabular}{lrrrr}
\toprule
Subset & $n$ & Accuracy & Mean confidence & Mean utility \\
\midrule
All images         & 25 & 0.96 & 3.44 & 3.48 \\
\midrule
Real (all)         & 15 & 1.00 & 3.73 & 3.80 \\
\quad Real (glioma)      &  5 & 1.00 & 2.80 & 2.60 \\
\quad Real (meningioma)  &  5 & 1.00 & 3.40 & 3.80 \\
\quad Real (pituitary)   &  5 & 1.00 & 5.00 & 5.00 \\
\midrule
Fake (all)         & 10 & 0.90 & 3.00 & 3.00 \\
\quad Fake (glioma)      &  3 & 1.00 & 2.33 & 2.33 \\
\quad Fake (meningioma)  &  2 & 0.50 & 2.00 & 2.00 \\
\quad Fake (pituitary)   &  5 & 1.00 & 3.80 & 3.80 \\
\bottomrule
\end{tabular}
\end{table}

\clearpage
\subsection{Evaluation of XAI Heatmaps}


The second part of the quantitative analysis focuses on the behavior of the explainability component, specifically the Grad-CAM heatmaps overlaid on real and synthetic MRI slices. For images presented with an XAI overlay, the neuroradiologist was asked not only to provide a tumor diagnosis, a confidence score, and a diagnostic utility rating, but also to judge whether the heatmap correctly highlighted the tumor region. This XAI rating was recorded using three categories: \emph{yes} (the heatmap clearly marked the tumor), \emph{partial} (the highlighted region overlapped the tumor only partially or ambiguously), and \emph{no} (the heatmap did not meaningfully localize the lesion).


\begin{table}[H]
\centering
\caption{Neuroradiologist performance by XAI pairing status and image type (all cases; unanswered counted as errors).}
\label{tab:hitl_xai_all_cases}
\begin{tabular}{lrrrr}
\toprule
Subset & $n$ & Accuracy & Mean confidence & Mean utility \\
\midrule
All images                 & 33 & 0.73 & 2.91 & 2.88 \\
\midrule
XAI paired (all)           & 18 & 0.78 & 3.28 & 3.33 \\
\quad real                 &  9 & 1.00 & 3.56 & 3.67 \\
\quad synthetic            &  9 & 0.56 & 3.00 & 3.00 \\
\midrule
XAI not paired (all)       & 15 & 0.67 & 2.47 & 2.33 \\
\quad real                 &  6 & 1.00 & 4.00 & 4.00 \\
\quad synthetic            &  9 & 0.44 & 1.44 & 1.22 \\
\bottomrule
\end{tabular}
\end{table}


\begin{table}[H]
\centering
\caption{Neuroradiologist performance by XAI pairing status and image type (answered cases only; rows with non-missing predicted tumor).}
\label{tab:hitl_xai_answered_only}
\begin{tabular}{lrrrr}
\toprule
Subset & $n$ & Accuracy & Mean confidence & Mean utility \\
\midrule
All images                 & 25 & 0.96 & 3.44 & 3.48 \\
\midrule
XAI paired (all)           & 14 & 1.00 & 3.93 & 4.00 \\
\quad real                 &  9 & 1.00 & 3.56 & 3.67 \\
\quad synthetic            &  5 & 1.00 & 4.60 & 4.60 \\
\midrule
XAI not paired (all)       & 11 & 0.91 & 2.82 & 2.82 \\
\quad real                 &  6 & 1.00 & 4.00 & 4.00 \\
\quad synthetic            &  5 & 0.80 & 1.40 & 1.40 \\
\bottomrule
\end{tabular}
\end{table}

\begin{table}[H]
\centering
\caption{Neuroradiologist performance by XAI heatmap rating and image type (all Grad-CAM cases; unanswered counted as errors). The XAI rating indicates whether the Grad-CAM heatmap correctly marked the tumor region. The rating ``yes'' is subdivided into ``full'' (heatmap fully covers tumor) and ``partial'' (heatmap partially overlaps tumor).}
\label{tab:hitl_xai_rating_real_fake_all_cases}
\begin{tabular}{lrrrr}
\toprule
Subset & $n$ & Accuracy & Mean confidence & Mean utility \\
\midrule
Real images           &  9 & 1.00 & 3.56 & 3.67 \\
\quad yes (overall)   &  6 & 1.00 & 3.50 & 3.67 \\
\quad\quad full       &  3 & 1.00 & 2.33 & 3.00 \\
\quad\quad partial    &  3 & 1.00 & 4.67 & 4.33 \\
\quad no              &  3 & 1.00 & 3.67 & 3.67 \\
\midrule
Synthetic images      &  9 & 0.56 & 3.00 & 3.00 \\
\quad yes (overall)   &  3 & 1.00 & 4.33 & 4.33 \\
\quad\quad full       &  2 & 1.00 & 4.00 & 4.00 \\
\quad\quad partial    &  1 & 1.00 & 5.00 & 5.00 \\
\quad no              &  6 & 0.33 & 2.33 & 2.33 \\
\bottomrule
\end{tabular}
\end{table}




The HITL evaluation further revealed that the clinical utility of the generated images is only partial. When neuroradiologist ratings are used as reference, approximately half of the synthetic images are judged to be diagnostically useful, while the remaining half are either ambiguous or clearly non diagnostic. Moreover, this utility is strongly class dependent: synthetic pituitary cases tend to be rated as realistic and clinically informative, whereas synthetic glioma and meningioma images receive substantially lower usefulness scores. These findings indicate that the generative model introduces a non-negligible class wise bias that must be carefully considered before integrating synthetic images into clinical pipelines.
